% ============================================================================
%  Chapter 4, second half: the replication in Python (RQ4).
%  Every number here is produced by
%     python vibebench/04_report/report_python.py
%     python vibebench/06_model/fit_mixed_model.py --lang py
%     python vibebench/04_report/py2_cross_language.py
%     python vibebench/labels/score_py2_labels.py --principle <p>
% ============================================================================

\section{Replication in a Second Language: Python (RQ4)}
\label{sec:python}

The Java results raise one question that Java alone cannot answer: whether they describe the
agents or describe how the agents write \emph{Java}. Change size dominates every principle, a
single agent effect survives adjustment, and risk sits inside single units; each of those could be
a property of a language in which classes are compulsory, contracts are declared, and a compiler
rejects a whole family of substitutability faults. This section applies the same framework, with
the same agents, the same curation policy, the same validation guards and the same statistical
model, to \textbf{Python}, and asks which of those results survive the move.

Python is the informative second language precisely because it removes what made Java convenient.
A Python module need not declare a class; a contract may be an \texttt{abc.ABC}, a
\texttt{typing.Protocol}, or nothing more than a base class whose methods raise
\texttt{NotImplementedError}; subtyping is checked by nobody; and every fault that
\texttt{javac} rejects at compile time is legal at run time. A result that holds in both languages
is therefore not an artefact of Java's type system.

\subsection{Corpus}
\label{sec:python-corpus}

The corpus is drawn from the same AIDev release with the same repository filters, at least 25
stars and a licence in the permissive allowlist of Section~\ref{sec:dataset}, keeping pull
requests that touch Python source by file extension. Two sampling choices differ from Java, and both
tighten the design rather than loosen it. The draw is balanced \emph{exactly} on agent and outcome,
1{,}400 merged, 700 rejected and 250 open requests per agent, 2{,}350 in all for each of the three;
and it is stratified by repository, so that no repository supplies more than 0.4\% of any
agent-outcome cell (Java's most concentrated cell took 17.2\% from one repository). The curated
corpus holds \textbf{7{,}050 pull requests over 3{,}974 repositories}, created between February
2025 and March 2026.

\begin{table}[htb]
\centering
\caption{Python corpus: analysable pull requests by agent and outcome.}
\label{tab:py-corpus}
\begin{tabular}{@{}lrrrr@{}}
\toprule
\textbf{Agent} & \textbf{Merged} & \textbf{Rejected} & \textbf{Open} & \textbf{Total} \\
\midrule
Claude Code    &   940 &   540 & 182 & 1{,}662 \\
GitHub Copilot &   885 &   407 & 178 & 1{,}470 \\
OpenAI Codex   & 1{,}069 &   549 & 196 & 1{,}814 \\
\midrule
\textbf{Total} & \textbf{2{,}894} & \textbf{1{,}496} & \textbf{556} & \textbf{4{,}946} \\
\bottomrule
\end{tabular}
\end{table}

Reconstruction uses the same four tiers as Java (Section~\ref{sec:dataset}). Of the 7{,}050 curated
requests, \textbf{4{,}946 over 2{,}921 repositories are analysable} (Table~\ref{tab:py-attrition}).
Unlike the Java corpus, the three agents are close to equally represented, and OpenAI Codex, the
smallest Java cell, is the largest Python one.

\begin{table}[htb]
\centering
\caption{Attrition from the curated Python corpus to the analysable one.}
\label{tab:py-attrition}
\begin{tabular}{@{}L{5.6cm}rL{5.2cm}@{}}
\toprule
\textbf{Status} & \textbf{PRs} & \textbf{Meaning} \\
\midrule
\texttt{ok} & \textbf{4{,}946} & analysable: valid before/after pair, Python touched \\
\texttt{no\_touched\_files} & 1{,}816 & no Python file left to analyse \\
\texttt{no\_touched\_modules} & 232 & files touched, none resolvable to a module \\
\texttt{no\_delta\_pair} & 40 & no tier could recover a distinct before tree \\
\texttt{repo\_too\_large} & 14 & tree exceeds the parsing budget \\
\texttt{archive\_failed} & 2 & a tree object missing, or a 4.9\,GB data tree refused \\
\midrule
\textbf{Total} & \textbf{7{,}050} & \\
\bottomrule
\end{tabular}
\end{table}

One component of that attrition is not a scope decision. \textbf{496 pull requests (7.0\%)} lost
their source entirely: the dataset carries no file list for them, they have no merge commit and no
recorded commits, and the branch that held them has been deleted from GitHub. The loss is uneven, 5.1\%
of merged requests against 11.0\% of rejected ones and 14.3\% of Copilot's requests against 3.6\% and
3.1\% for the other two agents, because a closed, unmerged branch is exactly what a repository
deletes. The corpus was deliberately \emph{not} topped up with replacements: replacements could only
be drawn from requests whose branches survived, which would make retrievability correlate with the
outcome the comparison measures.

\subsection{What Changes in the Instrument}
\label{sec:python-instrument}

The framework was built language-parameterised (Section~\ref{sec:language-availability}), and the
replication tests that claim. Only the fact extractor was replaced: Python source is parsed with the
standard library's \texttt{ast} module, with no JVM and no symbol solver, into the same record of
types, members, supertypes, imports and member accesses that the Java extractor emits. Every
detector shape, the delta harness, the attribution rules, the validation guards, the report and the
joint model are the Java code, unchanged.

Three things had to change, and each change follows from the language rather than from a choice.

\textbf{Thresholds were recalibrated, from Python's own distribution, before any label existed.}
A Java percentile is meaningless in Python: Python classes are smaller, and Python code expresses
dispatch differently. A calibration pass over all 7{,}050 requests recorded every metric with the
flags switched off; thresholds were set near the 90th percentile; the result tables were then
deleted and the pass re-run, so that no flag in this section predates its threshold.
Table~\ref{tab:py-thresholds} gives both sets.

\begin{table}[htb]
\centering
\caption{Detector thresholds in the two languages. Each set was calibrated on its own corpus before
labelling.}
\label{tab:py-thresholds}
\begin{tabular}{@{}lL{5.0cm}L{5.0cm}@{}}
\toprule
\textbf{Principle} & \textbf{Java} & \textbf{Python} \\
\midrule
SRP & rubric of seven size and cohesion metrics, $\geq$2 hits; class LOC 362, WMC 45, methods 20 &
same rubric, $\geq$2 hits; class LOC 250, WMC 36, methods 10, function CCN 12, parameters 7 \\
OCP & type chains $\geq$3 branches, tag chains $\geq$8 & type chains $\geq$3 branches, value chains
$\geq$5 \\
LSP & three signals: refused bequest, strengthened precondition, weakened postcondition & five
signals: refused bequest, widened return type, and three signature faults a compiler would reject \\
ISP & contract width $\geq$13 members & at least two of: more than 7 abstract members, mean
implementer coverage below 0.75, an implementer forced to stub \\
DIP & newly cyclic module pairs through a touched file & newly cyclic module pairs through a module
whose text changed \\
\bottomrule
\end{tabular}
\end{table}

\textbf{The Liskov detector regains the signals Java made unnecessary.} In Java a changed parameter
list is an overload and parameter types are invariant, so the Java detector does not look for
either (Section~\ref{sec:findings}, Finding~1). In Python both are live faults, so beside refused
bequest and a widened return type the Python detector adds three signals: a parameter removed,
a required parameter added, and a parameter type narrowed.

\textbf{A contract has to be recognised rather than read off a keyword.} Python has no
\texttt{interface}. The detector treats an \texttt{abc.ABC} subclass with abstract members and a
\texttt{typing.Protocol} as explicit contracts; a base class whose methods only raise
\texttt{NotImplementedError} is recorded as an implicit contract but excluded from the flag,
because in Python that shape is usually a template method with optional hooks. Unlike the Java flag,
the Python flag never fires on width alone: width must coincide with a second signal, low
implementer coverage or an implementer that stubs members. The Python and Java ISP rates are
therefore not the same quantity, and they are not compared as one.

\subsection{Instrument Validation}
\label{sec:python-validation}

The validation procedure is the Java one: a stratified sample of 100 flagged and 100 unflagged
entities per principle, judged by the same kind of reading judge against a written rubric, with
every verdict's reason recorded, every label traced to a batch file the judge was shown, the
detector's own answer withheld, and a mechanical audit run as a gate before judging. Two
principles admit a census of the flagged stratum. Interface Segregation flags only \textbf{80}
contracts in the entire Python corpus and all 80 were judged, so its precision is exact. For
Dependency Inversion, 100 of the 173 flagged pull requests were judged.
Table~\ref{tab:py-validation} gives the result over \textbf{980 judged entities}.

\begin{table}[htb]
\centering
\caption{Python detector agreement with the reading judge, beside the Java figures of
Table~\ref{tab:validation}. Recall is weighted by stratum size as in Java, with a 95\% Wilson
interval on the violation rate of the unflagged stratum. $^\dagger$ISP's flagged stratum is a census
of all 80 flagged contracts, so its precision carries no sampling error.}
\label{tab:py-validation}
\small
\begin{tabular}{@{}lrrrrll@{}}
\toprule
 & \multicolumn{4}{c}{\textbf{Python}} & & \textbf{Java} \\
\cmidrule(lr){2-5}
\textbf{Principle} & \textbf{judged} & \textbf{TP} & \textbf{FP} & \textbf{FN} &
\textbf{Precision / recall} & \textbf{Precision / recall} \\
\midrule
DIP & 200 & 100 & 0 & 0 & 1.000 / 1.000 \ (0.495--1.000) & 0.892 / 1.000 \\
LSP & 200 & 51 & 49 & 2 & 0.510 / 0.545 \ (0.255--0.813) & 0.420 / 0.352 \\
ISP$^\dagger$ & 180 & 34 & 46 & 0 & 0.425 / 1.000 \ (0.234--1.000) & 0.335 / 0.552 \\
OCP & 200 & 39 & 61 & 13 & 0.390 / 0.364 \ (0.262--0.490) & 0.330 / 0.733 \\
SRP & 200 & 11 & 89 & 0 & 0.110 / 1.000 \ (0.278--1.000) & 0.350 / 0.875 \\
\bottomrule
\end{tabular}
\end{table}

Four of the five detectors are at least as precise in Python as in Java. The exception is the one
that changed least.

\paragraph{Single Responsibility loses two thirds of its precision.} The rubric is the same, and so
are the kinds of classes it catches: large ones. What differs is what a large class \emph{is} in
Python. The 89 false positives are dominated by classes whose size is imposed from outside, a
\texttt{transformers} model head, a Home Assistant entity, a Lightning callback, a command class
whose methods are the subcommands of one tool; by clients for a single service or device, however
many commands the service exposes; and by configuration classes. Size and complexity measure the
same thing in both languages. Responsibility follows size less closely in Python, where frameworks
dictate the shape of the classes a project must write.

\paragraph{The Liskov detector confirms Finding 1 from the other side.} Of the 51 true positives,
\textbf{46 rest on a signal that \texttt{javac} would have rejected}: a parameter removed, a
required parameter added, or a parameter type narrowed. Only five rest on refused bequest alone. In
the language without the compiler, the faults the compiler removes are not merely possible; they
are most of what is actually there. The Java detector's three-signal design was not a compromise,
it was a measurement of what Java had already excluded.

\paragraph{Dependency Inversion is exact for a structural reason.} The single systematic false
positive found by the Java census was a cycle charged to a module declared as touched whose text
had not changed. The judging rule applies that check directly, comparing each charged module's text
between the two trees, and none of the 100 judged Python flags failed it: every one closes a new
import cycle through code the pull request changed, and no judged unflagged request does. The dominant shape is a new or edited module that imports its own
package root while the root re-exports it. The figure establishes that the cycle is real, not that
it is harmful, and the Python cycle findings below are read as ``introduced an import cycle''
throughout.

\paragraph{Precision by agent.} Table~\ref{tab:py-validation-agent} gives precision per agent. The
Open/Closed flags that carry the one agent effect surviving adjustment are \emph{most} often correct
for Claude Code, which rules out the reading that its effect is an artefact of noisier flags. The
Liskov column is uneven for a different reason: two Claude Code pull requests contribute 50 of the
100 flagged rows, one with 14 true positives in 30 and one with none in 20, and without those two
requests Liskov precision is 0.740.

\begin{table}[htb]
\centering
\caption{Python precision by agent, true positives over flags judged.}
\label{tab:py-validation-agent}
\begin{tabular}{@{}lccc@{}}
\toprule
\textbf{Principle} & \textbf{Claude Code} & \textbf{GitHub Copilot} & \textbf{OpenAI Codex} \\
\midrule
SRP & 0.130 \ (6/46)  & 0.077 \ (2/26)  & 0.107 \ (3/28) \\
OCP & 0.444 \ (20/45) & 0.368 \ (7/19)  & 0.333 \ (12/36) \\
LSP & 0.397 \ (29/73) & 0.733 \ (11/15) & 0.917 \ (11/12) \\
ISP & 0.479 \ (23/48) & 0.353 \ (6/17)  & 0.333 \ (5/15) \\
DIP & 1.000 \ (53/53) & 1.000 \ (16/16) & 1.000 \ (31/31) \\
\bottomrule
\end{tabular}
\end{table}

\subsection{Agent Differentiation under the Joint Model}
\label{sec:python-agents}

Table~\ref{tab:py-agent-rates} gives the raw flag rates. They look like the Java ordering, Claude
Code highest on four principles, and they are subject to the same warning: the agents write changes
of different size, 16.5 Python files per request on average for Claude Code against 7.7 for
Copilot and 9.6 for Codex.

\begin{table}[htb]
\centering
\caption{Python flag rate by agent. Pooled rate first, macro average over repositories in brackets.
SRP is reported over classes the agent authored, the other four over pull requests. A description,
not a ranking: read it only through Table~\ref{tab:py-model}.}
\label{tab:py-agent-rates}
\begin{tabular}{@{}lccc@{}}
\toprule
\textbf{Principle} & \textbf{Claude Code} & \textbf{GitHub Copilot} & \textbf{OpenAI Codex} \\
\midrule
SRP (authored classes) & 9.6\% (11.0\%) & \textbf{6.6\%} (10.3\%) & 8.9\% (8.1\%) \\
OCP & 9.7\% (9.6\%) & \textbf{3.5\%} (3.1\%) & 5.7\% (6.6\%) \\
ISP & 1.3\% (1.0\%) & 0.4\% (0.4\%) & \textbf{0.3\%} (0.1\%) \\
DIP & 5.8\% (5.6\%) & \textbf{2.0\%} (1.5\%) & 2.6\% (2.3\%) \\
LSP & 1.6\% (1.2\%) & \textbf{0.5\%} (0.3\%) & 0.8\% (0.8\%) \\
\bottomrule
\end{tabular}
\end{table}

The joint model of Section~\ref{sec:agent-differentiation} was fitted unchanged: a logistic
regression of the pull-request flag on agent, log files touched and months, estimated by
generalised estimating equations with an exchangeable correlation within repository and
cluster-robust standard errors, over all 4{,}946 requests in 2{,}921 repositories, with Copilot as
the reference agent. Python supports all five principles, including a pull-request-level Single
Responsibility flag and Liskov with 47 flagged requests (Table~\ref{tab:py-model}).

\begin{table}[htb]
\centering
\caption{Adjusted odds ratios from the joint model in Python, with 95\% confidence intervals. Each
agent coefficient is read against Copilot at the same change size, in the same period, in the same
repository. The last row repeats the Java estimate where one exists.}
\label{tab:py-model}
\small
\begin{tabular}{@{}lccccc@{}}
\toprule
\textbf{Term} & \textbf{SRP} & \textbf{OCP} & \textbf{ISP} & \textbf{DIP} & \textbf{LSP} \\
\midrule
Claude Code & 1.09 & \textbf{2.09} & 1.71 & \textbf{1.88} & 1.81 \\
 & (0.93--1.29) & (1.44--3.04) & (0.67--4.36) & (1.17--3.03) & (0.74--4.42) \\
OpenAI Codex & 1.06 & 1.46 & 0.53 & 1.12 & 1.38 \\
 & (0.90--1.25) & (0.98--2.16) & (0.16--1.80) & (0.67--1.86) & (0.54--3.52) \\
Size, per doubling & \textbf{2.06} & \textbf{2.01} & \textbf{2.17} & \textbf{2.04} & \textbf{2.18} \\
 & (1.94--2.18) & (1.88--2.15) & (1.82--2.58) & (1.89--2.21) & (1.89--2.51) \\
Era, per month & \textbf{1.05} & 0.99 & 0.94 & 1.06 & 0.96 \\
\midrule
flagged / total & 2{,}532 / 4{,}946 & 315 / 4{,}946 & 32 / 4{,}946 & 173 / 4{,}946 & 47 / 4{,}946 \\
\midrule
Java, Claude Code & n/a & 2.22 & 1.11 & 1.29 & n/a \\
Java, per doubling & n/a & 1.89 & 1.88 & 1.95 & n/a \\
\bottomrule
\end{tabular}
\end{table}

\textbf{Change size dominates again, at almost the same magnitude.} Each doubling of the files a
pull request touches doubles the odds of a flag, 2.01 to 2.18 on all five principles
($p < 10^{-4}$ throughout), against 1.88 to 1.95 in Java. It is the largest effect in either
language, and it is now a result about agent-authored pull requests rather than about Java.

\textbf{The Open/Closed agent effect reproduces.} Claude Code's Open/Closed odds are 2.09 times
Copilot's at the same size, date and repository (1.44--3.04, $p = 0.0001$), against 2.22
(1.21--4.07) in Java. The two intervals overlap almost entirely, and the Python interval is narrower.
Of the agent differences Table~\ref{tab:py-agent-rates} suggests, this is the one that is a property
of the code rather than of how much code was written, in both languages.

\textbf{One difference appears in Python that Java did not show.} Claude Code's odds of introducing
an import cycle are 1.88 times Copilot's (1.17--3.03, $p = 0.010$); in Java the corresponding ratio
was 1.29 with an interval straddling 1. Every other agent coefficient in Python, for every agent
and principle, has an interval that includes 1.

\textbf{Era is inert for the four delta principles, not for Single Responsibility.} The monthly
coefficient is indistinguishable from 1 for OCP, ISP, DIP and LSP, as it was in Java. The
pull-request-level Single Responsibility flag, which Java did not model, rises by 5\% per month
($p < 10^{-4}$) over a thirteen-month window.

\subsection{Merged Code against Rejected Code}
\label{sec:python-outcome}

The outcome comparison is made as in Java: on the tier-matched subset, where merged and not-merged
requests recover their \emph{before} tree the same way, and by stratum
(Table~\ref{tab:py-outcome}).

\begin{table}[htb]
\centering
\caption{Python flag rate by outcome stratum on the tier-matched subset (tiers 2, 3 and 4), with
flagged counts in brackets. SRP is a per-class measure and is not split this way.}
\label{tab:py-outcome}
\begin{tabular}{@{}lrcccc@{}}
\toprule
\textbf{Stratum} & \textbf{n} & \textbf{OCP} & \textbf{ISP} & \textbf{DIP} & \textbf{LSP} \\
\midrule
\textbf{merged} & 1{,}734 & 4.9\% (85) & 0.4\% (7) & 3.1\% (54) & 0.5\% (8) \\
\midrule
\texttt{changes\_requested} & 33 & 9.1\% (3) & 0.0\% (0) & 9.1\% (3) & 0.0\% (0) \\
\texttt{reviewed\_then\_closed} & 163 & 7.4\% (12) & 1.8\% (3) & 8.6\% (14) & 1.2\% (2) \\
\texttt{closed\_unreviewed} & 1{,}300 & 9.1\% (118) & 0.9\% (12) & 4.5\% (59) & 1.7\% (22) \\
\midrule
\texttt{open} & 553 & 10.8\% (60) & 1.1\% (6) & 5.2\% (29) & 2.2\% (12) \\
\bottomrule
\end{tabular}
\end{table}

\textbf{Here the two languages disagree.} In Java, merged code was indistinguishable from code closed
without review. In Python it is not: merged pull requests carry a lower flag rate than the
unreviewed, reviewed-then-closed and open strata on every one of the four principles, and against
the 1{,}300 unreviewed closures the Open/Closed rate is 4.9\% against 9.1\%. Only the 33
\texttt{changes\_requested} requests, too few to rank, sit lower on two principles.

Part of that gap is change size, since rejected requests are larger (median four files against
three). Adding size and agent to the comparison, clustered by repository as before, leaves
rejected requests with higher Open/Closed odds (odds ratio 1.57, 1.15--2.15, $p = 0.004$) and higher
Liskov odds (2.43, 1.08--5.46, $p = 0.03$), while the Interface Segregation (1.71, 0.75--3.91) and
Dependency Inversion (1.36, 0.92--2.03) differences are not distinguishable from none. The
precision of the Open/Closed flag also differs by outcome, 0.239 on merged requests against 0.500 on
rejected ones, and it differs in the direction that widens the gap once flags are corrected to
genuine violations. Because most not-merged requests were never reviewed, the result says that
merged Python code is structurally cleaner than the code that was abandoned; it does not say that
reviewers rejected code for its structure.

\subsection{Locus of Risk}
\label{sec:python-locus}

Table~\ref{tab:py-present-introduced} repeats the present-versus-introduced contrast of
Section~\ref{sec:locus-of-risk}, correcting both columns by the Python precision of
Table~\ref{tab:py-validation}.

\begin{table}[htb]
\centering
\caption{Python violations present in touched code against pull requests that introduced one,
corrected for measured precision, beside the Java values of Table~\ref{tab:present-introduced}.
Present entities are distinct entities. For SRP, ``introduced'' means a flagged class the pull
request newly authored, which the Java table did not report.}
\label{tab:py-present-introduced}
\small
\begin{tabular}{@{}lrrrr@{}}
\toprule
 & \multicolumn{3}{c}{\textbf{Python}} & \textbf{Java} \\
\cmidrule(lr){2-4}
\textbf{Principle} & \textbf{Flagged present} & \textbf{Est.\ genuine} & \textbf{PRs introduced} &
\textbf{PRs introduced} \\
\midrule
SRP & 8{,}208 classes & $\approx$900 & 449 (1.0\% corrected) & n/a \\
OCP & 2{,}580 chains & $\approx$1{,}000 & 315 (2.5\% corrected) & 2.2\% \\
ISP & 80 contracts & $\approx$34 & 32 (0.3\% corrected) & 2.1\% \\
LSP & 2{,}069 override pairs & $\approx$1{,}050 & 47 (0.5\% corrected) & 0.4\% \\
DIP & 173 cycles & 173 & 173 (3.5\% corrected) & 4.1\% \\
\bottomrule
\end{tabular}
\end{table}

The contrast holds, and for the relation-level principles it holds at almost the same ratio.
\textbf{194 refused-bequest relationships are present in the Python code these agents touched, and
five pull requests introduced one}, against 681 and 16 in Java: roughly forty present for every one
introduced, in both languages. Across all five principles, between 0.3 and 3.5 per cent of pull
requests introduce a genuine violation into code that already holds tens to thousands.

The authored-versus-edited comparison also points the same way. Classes the agents authored are
flagged at 8.9\% against 12.1\% for the pre-existing classes they edited, and at 9.7\% against 23.6\%
under repository-mix control. The gap is narrower than Java's (7.7\% against 13.5\% pooled, 5.6\%
against 20.9\% under mix control), but it has the same sign on both measures.

\subsection{The Refinement Lesson, Tested out of Language}
\label{sec:python-refinements}

Section~\ref{sec:instrument-validation} found that a categorical restriction (Liskov restricted to
refused bequest) survived a holdout while a numeric one (Interface Segregation restricted to 30 or
more members) did not. Both rules were fixed on Java data before any Python label existed. Applying
them, unchanged, to the Python judgements is therefore a holdout in a second language, and no part
of it can have been fitted (Table~\ref{tab:py-refinements}).

\begin{table}[htb]
\centering
\caption{The two Java refinements applied without change to the Python judgements. Java counts all
contract members; Python counts abstract members, the width its detector records.}
\label{tab:py-refinements}
\begin{tabular}{@{}lll@{}}
\toprule
\textbf{Refinement} & \textbf{Java (fitted $\rightarrow$ holdout)} & \textbf{Python} \\
\midrule
LSP, restricted to refused bequest & 1.000 $\rightarrow$ 1.000 & \textbf{0.857} \ (6/7) \\
ISP, restricted to $\geq$30 members & 0.862 $\rightarrow$ 0.704 & \textbf{0.200} \ (1/5) \\
\bottomrule
\end{tabular}
\end{table}

The categorical rule carries over, at six correct in seven; the one exception is an override that
repeats a refusal its base method already makes for the same case. The numeric threshold does not carry over
at all. Both Python counts are small, and neither is offered as a new estimate. What they add is
direction: a rule that names a shape of code transferred to a language with a different type system
and different contract idioms, and a rule that names a quantity did not survive even within one.

\subsection{What Generalises}
\label{sec:python-summary}

Table~\ref{tab:cross-language} collects the comparison.

\begin{table}[htb]
\centering
\caption{The Java results tested in Python.}
\label{tab:cross-language}
\begin{tabular}{@{}L{4.2cm}L{4.4cm}L{4.4cm}@{}}
\toprule
\textbf{Result} & \textbf{Java} & \textbf{Python} \\
\midrule
Change size dominates & OR 1.88--1.95 per doubling & \textbf{holds}: OR 2.01--2.18 \\
Open/Closed agent effect & Claude Code OR 2.22 (1.21--4.07) & \textbf{holds}: 2.09 (1.44--3.04) \\
Other agent effects & none survive adjustment & DIP survives (1.88, 1.17--3.03) \\
Risk present, not introduced & 681 present / 16 introduced & \textbf{holds}: 194 / 5 \\
Authored classes flagged less & 7.7\% vs 13.5\% & \textbf{holds}: 8.9\% vs 12.1\% \\
Categorical rule beats numeric & 1.000 vs 0.704 on holdout & \textbf{holds}: 0.857 vs 0.200 \\
Compiler removes LSP faults & inferred from Java's signals & \textbf{confirmed}: 46 of 51 true
positives are compiler-rejectable \\
Merged equals unreviewed & indistinguishable & \textbf{does not hold}: merged lower \\
Detector precision & 0.330--0.892 & 0.110--1.000; SRP drops to 0.110 \\
\bottomrule
\end{tabular}
\end{table}

The results about change size, the Open/Closed effect, the locus of risk and the fate of fitted
rules are properties of agent-written code, not of Java. The result about review is not stable
across languages, and the Single Responsibility screen transfers badly. Section~\ref{sec:rq-answers}
draws the conclusions.
